\chapter{Model-free multi-currency portfolio returns} 
\label{chap:model-free-returns}

In this chapter, we derive a model-free representation of one-period portfolio returns for an investor who holds assets denominated in multiple currencies and can take currency forward positions for hedging.\ The derivations follow the methodology presented in \cite{BurkhardtUlrych2023} and \cite{LucescuUlrych2026}, using simple returns and assume that rebalancing occurs at the end of each period. The representation is used in subsequent chapters to formulate and analyze the results of the multi-currency portfolio optimization problem.

\section{Unhedged Portfolio Returns}

We consider an international investor that allocates his wealth across $N$ assets denominated in up to $M+1$ currencies.\ We fix a discrete time index $t\in\{0,1,2,\dots\}$ and one investment period from $t$ to $t{+}1$. Let the investor's \emph{domestic currency} be indexed by $0$ and let the rest of the $M$ foreign currencies be indexed by $\mathcal{C} \coloneqq \{1,2,\dots,M\}$.

Let $P_{i,t}>0$ be the price of asset $i$ at time $t$ in its local currency $c_i \in \mathcal{C}\cup \{0\}$. Define the corresponding \emph{simple local currency asset return} over $(t,t{+}1]$ as
\begin{equation*}
R^{\text{LC}}_{i,t+1} \coloneqq \frac{P_{i,t+1}-P_{i,t}}{P_{i,t}}.
\end{equation*}
Thus $P_{i,t+1}=P_{i,t}(1+R^{\text{LC}}_{i,t+1})$.

For each foreign currency $c\in\mathcal{C}$, let $S_{c,t} > 0$ denote the exchange rate at time $t$, quoted as units of domestic currency per one unit of foreign currency $c$. The associated \emph{simple currency return} over $(t,t{+}1]$ is
\begin{equation*}
e_{c,t+1} \coloneqq \frac{S_{c,t+1}-S_{c,t}}{S_{c,t}}.
\end{equation*}
Equivalently, $S_{c,t+1} = S_{c,t}(1+e_{c,t+1})$. By definition for the domestic currency, we set $S_{0,t}\equiv 1$ and $e_{0,t+1}\equiv 0$.

The domestic currency value of holding one unit of asset $i$ at time $t$ is $P_{i,t}\,S_{c_i,t}$. Therefore, the \emph{unhedged domestic currency simple return} of asset $i$ is
\begin{align}
R^{u}_{i,t+1}
&=
\frac{P_{i,t+1}S_{c_i,t+1} - P_{i,t}S_{c_i,t}}{P_{i,t}S_{c_i,t}}
\nonumber\\
&=
\frac{P_{i,t}(1+R^{\text{LC}}_{i,t+1}) S_{c_i,t}(1+e_{c_i,t+1}) - P_{i,t}S_{c_i,t}}{P_{i,t}S_{c_i,t}}
\nonumber\\
&=
R^{\text{LC}}_{i,t+1} + e_{c_i,t+1} + R^{\text{LC}}_{i,t+1}\,e_{c_i,t+1}.
\label{eq:unhedged-asset-return}
\end{align}

Let $\bm{x}_t\in\R^N$ be the vector of asset portfolio weights chosen at time $t$ and held until $t{+}1$, where $x_{i,t}$ is the fraction of total portfolio value invested in asset $i$ at time $t$. We impose the standard budget constraint $\bm{1}^\T \bm{x}_t = 1$.

Assuming no rebalancing between times $t$ and $t{+}1$, the domestic currency portfolio's unhedged return is given by
\begin{equation*}
R^{u}_{\mathcal{P},t+1} = \sum_{i=1}^N x_{i,t}\,R^{u}_{i,t+1}.
\end{equation*}

\section{Hedged Portfolio Returns}

To manage currency risk, the investor can take positions in one-period currency forward contracts. A forward contract on foreign currency $c$ entered at time $t$ settles at time $t{+}1$ with payoff (in domestic currency) of $F_{c,t} - S_{c,t+1}$, where $F_{c,t}$ is the forward rate quoted at time $t$. The forward rate is typically set such that the contract has zero value at inception, i.e.\ $F_{c,t} = S_{c,t}(1+r_t)/(1+r_{c,t})$, where $r_t$ and $r_{c,t}$ are the domestic and foreign risk-free rates, respectively. The \emph{excess currency return} over the period is $e_{c,t+1} - f_{c,t}$, where $f_{c,t} = (F_{c,t} - S_{c,t})/S_{c,t}$ is the forward premium. The investor can choose the forward position $\tau_{c,t}$ (expressed as a fraction of total portfolio value) for each foreign currency $c$ to hedge against currency risk. The \emph{hedged} portfolio return is then the sum of the unhedged asset returns and the payoffs from the forward positions:
\begin{equation*}
R^{h}_{\mathcal{P},t+1} = \sum_{i=1}^N x_{i,t}\,R^{u}_{i,t+1} + \sum_{c=1}^M \tau_{c,t}\,(f_{c,t} - e_{c,t+1}).
\end{equation*}
For each foreign currency $c\in\mathcal{C}$, assume that at time $t$ the investor can enter a one-period currency forward contract that settles at $t{+}1$, with forward rate $F_{c,t}$ quoted in domestic currency per unit of foreign currency $c$. Define the \emph{forward premium} (simple rate) by
\begin{equation*}
f_{c,t} \coloneqq \frac{F_{c,t}-S_{c,t}}{S_{c,t}}.
\end{equation*}
Let $\bm{\tau}_t\in\R^M$ denote forward positions expressed as fractions of total portfolio value, where $\tau_{c,t}$ represents the relative notional of the \emph{short} forward in currency $c$ at time $t$. With this sign convention, $\tau_{c,t}>0$ represents hedging (shorting foreign currency forward). The \emph{hedged} portfolio return is
\begin{equation*}
R^{h}_{P,t+1}
=
\sum_{i=1}^N x_{i,t}\,R^{u}_{i,t+1}
+
\sum_{c=1}^M \tau_{c,t}\,(f_{c,t}-e_{c,t+1}).
\end{equation*}
This representation is model-free: it holds pathwise for realized returns and does not assume any distributional structure.

\section{Matrix Representation and Decomposition}

To express the returns in a compact form suitable for optimization, we define the currency incidence matrix $\bm{C}\in\{0,1\}^{N\times M}$ by
\begin{equation*}
C_{i,c} \coloneqq
\begin{cases}
1, & \text{if asset $i$ is denominated in foreign currency $c$},\\
0, & \text{otherwise}.
\end{cases}
\end{equation*}
The portfolio's \emph{implicit foreign-currency exposure} is $\bm{w}_t \coloneqq \bm{C}^\T \bm{x}_t \in \R^M$, where $w_{c,t} = \sum_{i:\,c_i=c} x_{i,t}$.

Collect returns into vectors $\bm{R}^{\text{LC}}_{t+1} \in \R^N$, $\bm{e}_{t+1} \in \R^M$, and the forward premiums $\bm{f}_t \in \R^M$. Define the currency-return vector aligned to assets via $\bm{C}\bm{e}_{t+1}\in\R^N$, where the $i$th component is $e_{c_i,t+1}$ for foreign assets and $0$ for domestic assets. Using \eqref{eq:unhedged-asset-return}, the vector of \emph{unhedged domestic returns} across assets is
\begin{equation*}
\bm{R}^{u}_{t+1}
=
\bm{R}^{\text{LC}}_{t+1}
+
\bm{C}\bm{e}_{t+1}
+
(\bm{R}^{\text{LC}}_{t+1}\had\bm{C}\bm{e}_{t+1}),
\end{equation*}
where $\had$ denotes the Hadamard product. 

